% ============================================
% DOC-01 (EN)
% English translation — SPEC-935-R355
% ============================================
\section*{\textbf{DOC-01} --- The Exhumation}
\markboth{DOC-01 --- The Exhumation}{}

\noindent \textit{Notes of gravedigger Genésio Barbosa}

\noindent \textbf{Notes of gravedigger Genésio Barbosa}
\textit{Literal transcription of the statement given on June 14, 1979 to the administration of the Hospital-Colony. Document filed under code C-1260-79.}

\vspace{0.5em}\noindent\rule{\textwidth}{0.5pt}\vspace{0.5em}

\noindent It was the third grave of the day. You can tell the old ones because the soil has settled. Recent graves, the earth stays loose, you can almost dig with your hands. Not the old ones. The earth hardens, forms a crust on top, and when you sink the shovel in it feels like you are going to hit stone.

\noindent The 1,260 one was old. Very old. According to the map, 1917. One of the first buried there.

\noindent I dug. I dug about forty centimeters and the shovel hit wood. I dug more. The coffin --- it was a coffin, but it was not. You know when wood rots and becomes earth? It does not disappear. It turns into something else. A dark paste that still has the shape of what it was.

\noindent I opened it with the shovel. Not out of disrespect. It is that you could not pull the lid off --- it was already part of the ground.

\noindent What I saw inside was not a skeleton.

\noindent It was a body.

\noindent No, I am not crazy. It was a body. The skin was dark, almost black, but it was \textit{whole}. As if he had been buried last week. You can ask Mr. Nilo, who was with me. He saw it too. The skin of the face still covered the bones. The mouth was open. The eyes ---

\noindent The eyes were not there. The sockets were empty. But inside them, at the bottom, there was a glow. As if someone had lit a candle in there and blown it out. No, it was not a candle. It was a reflection. Yellow. Very faint.

\noindent And on the chest of the body --- tied with a strip of rotten leather, pressed tight against the bone --- there was a notebook. The leather crossed the ribs, under the jacket, as if the man had been buried with the book glued to his chest. I had to cut the strap with a knife. The leather crumbled in my hand.

\noindent I know it sounds like a lie. A notebook buried with a patient in 1917, tied to the body, sixty-two years later, with the pages still \textit{moist}? I would not believe it. But you are holding the notebook. You saw the pages. The ink.

\noindent You know I am not lying.

\noindent I picked up the notebook. When I touched it, I felt a warmth. It was not the warmth of the sun. It was a damp warmth, like that of a covered wound. As if the notebook were alive in there. Growing.

\noindent I opened the first page.

\noindent The handwriting was a child's. Crooked, rising and falling, like someone who learned to write on his own. A signature at the top:

\noindent \textit{Joaquim Antônio Correia --- Patient 1,260}

\noindent And a date.

\noindent \textit{1917 --- the year I came here}

\noindent But the diary went on after that. Page after page. The last entry was from 1979.

\noindent \textit{Three days before I dug him up.}

\noindent You understand? The man had been dead for 62 years. And the diary kept being written. In the grave. In the dark. With the fingers of a skeleton.

\noindent I did not sleep that night.

\noindent I have not slept since.

\noindent You can keep the diary. I do not want it near me anymore.

\noindent --- \textit{Genésio Barbosa, 48 years old, gravedigger at the Barbacena Hospital-Colony for 22 years.}

\vspace{0.5em}
\noindent\textit{\textbf{↪ Routes:} \rota{MEM-04} \rotasep \rota{LUC-02}}
